\documentclass[11pt,a4paper]{article}

%% -------- Core Packages --------
\usepackage{fontspec}
\usepackage{polyglossia}
\usepackage{amsmath,amssymb,amsthm}
\usepackage{geometry}
\usepackage{graphicx}
\usepackage{xcolor}
\usepackage{titlesec}
\usepackage{fancyhdr}
\usepackage{enumitem}
\usepackage{array,booktabs,longtable,multirow,colortbl}
\usepackage{hyperref}
\usepackage{microtype}
\usepackage{tocloft}
\usepackage{indentfirst}
\usepackage{xeCJK}

%% -------- Page Geometry --------
\geometry{
  paperwidth=210mm,paperheight=297mm,
  left=22mm,right=22mm,top=28mm,bottom=28mm,
  headheight=14pt,footskip=30pt
}

%% -------- Language Setup --------
\setmainlanguage{english}

%% -------- Font Configuration --------
\setmainfont[
  UprightFont=* Regular,
  BoldFont=* Bold,
  ItalicFont=* Italic,
  BoldItalicFont=* Bold Italic
]{Noto Serif}
\setsansfont[
  UprightFont=* Regular,
  BoldFont=* Bold,
  ItalicFont=* Italic
]{Noto Sans}
\setmonofont{DejaVu Sans Mono}

%% CJK Support via xeCJK -- use specified font path
\setCJKmainfont[AutoFakeBold=true,AutoFakeSlant=true]{Noto Sans CJK SC}
\setCJKsansfont[AutoFakeBold=true,AutoFakeSlant=true]{Noto Sans CJK SC}
\setCJKmonofont[AutoFakeBold=true,AutoFakeSlant=true]{Noto Sans CJK SC}

%% -------- Colors --------
\definecolor{titlecolor}{RGB}{45,55,72}
\definecolor{sectioncolor}{RGB}{55,65,81}
\definecolor{notecolor}{RGB}{120,53,15}
\definecolor{rulecolor}{RGB}{180,160,140}
\definecolor{volcolor}{RGB}{80,40,20}
\definecolor{probbg}{RGB}{250,248,245}

%% -------- Section Formatting --------
\titleformat{\section}
  {\normalfont\Large\bfseries\color{volcolor}}
  {\thesection}{1em}{}
\titleformat{\subsection}
  {\normalfont\large\bfseries\color{sectioncolor}}
  {\thesubsection}{1em}{}
\titleformat{\subsubsection}
  {\normalfont\normalsize\bfseries\color{sectioncolor}}
  {\thesubsubsection}{1em}{}

%% -------- Header/Footer --------
\pagestyle{fancy}
\fancyhf{}
\fancyhead[L]{\small\textsc{Ceyuan Haijing}}
\fancyhead[R]{\small\thepage}
\renewcommand{\headrulewidth}{0.4pt}
\renewcommand{\footrulewidth}{0pt}

%% -------- Custom Commands --------
\newcommand{\longline}{\noindent\rule{\textwidth}{0.4pt}}
\newcommand{\shortline}{\noindent\rule{0.5\textwidth}{0.4pt}\par}
\newcommand{\cnote}[1]{\textcolor{notecolor}{\textit{#1}}}

%% Problem/solution environment for Chinese mathematical texts
\newcounter{probnum}
\newcommand{\chnum}[1]{%
  \ifcase#1 零\or 一\or 二\or 三\or 四\or 五\or 六\or 七\or 八\or 九\or 十\or
    十一\or 十二\or 十三\or 十四\or 十五\or 十六\or 十七\or 十八\or 十九\or 二十\or
    二十一\or 二十二\or 二十三\or 二十四\or 二十五\or 二十六\or 二十七\or 二十八\or 二十九\or 三十\or
    三十一\or 三十二\or 三十三\or 三十四\or 三十五\fi}

\newcommand{\problem}[1]{\stepcounter{probnum}\par\noindent\textbf{第\chnum{\value{probnum}}問：}#1\par\vspace{0.2cm}}
\newcommand{\answer}[1]{\par\noindent\textbf{答曰：}#1\par\vspace{0.2cm}}
\newcommand{\working}[1]{\par\noindent\textbf{草曰：}#1\par\vspace{0.2cm}}
\newcommand{\method}[1]{\par\noindent\textbf{法曰：}#1\par\vspace{0.3cm}}

%% Polynomial display helper (tian yuan notation)
\newcommand{\tianyuan}[2]{\raisebox{-0.6ex}{\scriptsize$\begin{array}{c}#1\\\hline #2\end{array}$}}
\newcommand{\ty}[1]{\tianyuan{#1}{1}}
\newcommand{\tym}[2]{\tianyuan{#1}{#2}}

%% -------- Hyperref Setup --------
\hypersetup{
  colorlinks=true,
  linkcolor=sectioncolor,
  citecolor=sectioncolor,
  urlcolor=sectioncolor,
  pdfencoding=auto
}

%% -------- TOC Depth --------
\setcounter{tocdepth}{3}
\setcounter{secnumdepth}{3}

%% ===================================================================

\newcommand{\vs}{\vspace{0.5em}}
\begin{document}

%% ===================================================================
%% TITLE PAGE
%% ===================================================================
\begin{titlepage}
\begin{center}
\vspace*{2cm}

{\fontsize{14}{18}\selectfont 欽定四庫全書}

\vspace{1.5cm}

{\fontsize{12}{16}\selectfont 子部}

\vspace{2cm}

\rule{0.8\textwidth}{1.5pt}

\vspace{1cm}

{\fontsize{28}{36}\selectfont\bfseries\color{titlecolor} 測圓海鏡}

\vspace{0.5cm}

{\fontsize{16}{22}\selectfont 卷四至卷六}

\vspace{0.3cm}

{\fontsize{12}{16}\selectfont\color{notecolor} （底勾、大股、大勾三卷全編）}

\vspace{1cm}

\rule{0.8\textwidth}{1.5pt}

\vspace{2cm}

{\fontsize{14}{18}\selectfont 元\quad 李冶\quad 撰}

\vspace{0.5cm}

{\fontsize{11}{15}\selectfont\color{notecolor} 字仁卿，號敬齋，眞定欒城人}

\vspace{3cm}

{\fontsize{11}{15}\selectfont\color{notecolor}
\begin{tabular}{rl}
詳校官 & 欽天監博士 臣 張天樞 \\
靈臺郎 臣 倪廷梅 & 覆勘 \\
總校官 知縣 臣 繆琪 & \\
校對官 五官靈臺郎 臣 陳際新 & \\
謄錄監生 臣 施華 & \\
\end{tabular}
}

\vfill

{\small 依《欽定四庫全書》本}

\vspace{0.3cm}

{\small\color{notecolor} 卷四：第五十一問至第六十七問}

{\small\color{notecolor} 卷五：第六十八問至第八十五問}

{\small\color{notecolor} 卷六：第八十六問至第一百三問}

\end{center}
\end{titlepage}

%% ===================================================================
%% TABLE OF CONTENTS
%% ===================================================================
\tableofcontents
\newpage

%% ===================================================================
%% INTRODUCTION
%% ===================================================================
\section*{引言}
\addcontentsline{toc}{section}{引言}

《測圓海鏡》十二卷，元李冶撰。冶字仁卿，號敬齋，眞定欒城人。
金末進士，入元官翰林學士。此書乃其所著天元術之傑作，以勾股測圓
之法，設問答以明天元一之術。全書凡一百七十問，皆設城池圓形，以
人物行步之數推求圓徑。

本編收錄卷四、卷五、卷六三卷，凡五十三問：
\begin{itemize}
  \item \textbf{卷四}：底勾一十七問（第五十一問至第六十七問）
  \item \textbf{卷五}：大股一十八問（第六十八問至第八十五問）
  \item \textbf{卷六}：大勾一十八問（第八十六問至第一百三問）
\end{itemize}

書中設問之例：假令有圓城一座，甲乙二人各從城門出行，或東或南，
行若干步，遙相望見，或斜行相會。只知若干步數，餘皆不知，以求
圓城之徑。每問皆有\textbf{問}、\textbf{答}、\textbf{草}、\textbf{法}四部分：
\begin{itemize}
  \item \textbf{問}：設問條件，敘述甲乙二人行步之數及相會之狀
  \item \textbf{答曰}：所得答案，即圓城之徑
  \item \textbf{草曰}：以天元術列方程之詳細演算，為全書精華所在
  \item \textbf{法曰}：總結算法之簡要公式，便於後學記誦
\end{itemize}

\vspace{0.3cm}
\noindent
\textbf{天元術記法說明：}書中以籌算式記多項式，自下而上排列，
常數項在下，一次項在上，二次又在上，依此類推。如
\tianyuan{480}{1}表示 $480 - x$，\tianyuan{0\quad0\quad48000}{1\quad0\quad0}
表示 $48000 - x^2$ 之類。所謂「立天元一為某數」，即設未知數 $x$ 為
所求之量。

\vspace{0.5cm}
\longline

\newpage

%% Reset problem counter for Volume 4 (starts at problem 51 overall)
\setcounter{probnum}{0}

%% ===================================================================
%% VOLUME 4
%% ===================================================================
\section{測圓海鏡卷四}
\label{sec:vol4}

\begin{center}
{\Large\bfseries\color{volcolor} 底勾一十七問}

\vspace{0.2cm}

{\small\color{notecolor} （第五十一問至第六十七問）}
\end{center}

\vspace{0.3cm}
\longline
\vspace{0.3cm}

\noindent
\textbf{【題意總說】}卷四論「底勾」之術。所謂底勾，乃勾股形之最下
勾股，與圓徑相切之基本勾股形也。凡一十七問，皆以底勾為基，求
圓城之徑。底勾者，通勾股形之勾邊，亦即大勾，為諸勾股形之根本。
諸問或知甲東行之數，或知乙南行之數，或知斜行之數，三者之中知其二，
即可立天元一，建方程而開方得圓徑。

\vspace{0.3cm}

%% --- Problem 51 (Vol 4, Problem 1) ---
\problem{
或問：乙出南門東行不知步數而立，甲出北門東行二百步見之，就乙斜行
二百七十二步與乙相會。問圓城之徑幾何？
}

\answer{城徑二百四十步。}

\working{
識別得二行相減餘，即乙出南門東行數也。以甲東行二百步減於就乙斜行
二百七十二步，餘七十二步，為乙東行數。以乘甲東行步，得
$72 \times 200 ={}$一萬四千四百步，又四之，得五萬七千六百步為實。
以平方開之，得 $\sqrt{57600} ={}$二百四十步，即城徑也。

\textbf{驗算：}城徑二百四十步，半徑一百二十步。乙東行七十二步，
甲東行二百步，甲比乙多行 $200 - 72 = 128$ 步。以勾股術驗之，
斜行為弦，$\sqrt{128^2 + 120^2} = \sqrt{16384 + 14400} = \sqrt{30784}$，
復以半徑加之，合得斜行步數，其術驗矣。
}

\method{
二行差數乘甲東行又四之，為平方實，開方得全徑。術曰：以就乙斜行步
內減甲東行步，餘為乙東行數；以此數乘甲東行步為實，四之；開平方，
即得城徑。
}

%% --- Problem 52 (Vol 4, Problem 2) ---
\problem{
或問：乙出南門東行一百二十步而立，甲出北門東行二百步見之。問城徑。
}

\answer{城徑二百四十步。}

\working{
立天元一為城徑，以減於二之甲東行步，得\tianyuan{480}{1}為兩个小差。
以乙東行一百二十步乘之，得\tianyuan{57600}{1}為城徑冪，寄左。
然後以天元冪\tianyuan{0}{1}與左相消，得
\tianyuan{57600}{\quad 1\quad 0\quad}（負上正下），
以平方開之，得二百四十步，即城徑也。

\textbf{釋義：}二之甲東行者，以甲東行為勾，其倍即通勾與圓徑之差。
以乙東行乘之者，乙東行即小勾，乘大差得大勾之冪，亦即城徑冪之半。
故以天元冪相消，開平方即得。
}

\method{
半之乙東行步乘甲東行為實，半乙東行為從，一步常法，得半徑。
術曰：置乙東行步折半，以乘甲東行步為實；以半乙東行步為從法；
一步為隅法；開平方得半徑，倍之即城徑。
}

%% --- Problem 53 (Vol 4, Problem 3) ---
\problem{
或問：乙從坤隅南行三百六十步，甲出北門東行二百步見之。問城徑幾何？
}

\answer{城徑二百四十步。}

\working{
立天元一為半徑，減於乙南行三百六十步之半，得\tianyuan{180}{1}為小差。
半乙南行步，得一百八十步，以乘小差，得\tianyuan{32400}{180}為半徑冪，
寄左。然後以天元冪\tianyuan{0}{1}與左相消，得
\tianyuan{32400}{180\quad 1\quad 0}（上負下正），
開平方得一百二十步，倍之得全徑二百四十步。

\textbf{釋義：}此問以乙南行為股，甲東行為勾。半乙南行即大股之半，
去半徑餘為小差。以小差乘半股得半徑冪，此天元術之妙用也。
}

\method{
兩行步相乘為實，甲東行為從，乙南行為隅，開平方得半徑。
術曰：以甲東行步乘乙南行步為實，以甲東行步為從，乙南行步為隅法；
開平方得一百二十步，即半徑；倍之即城徑。
}

%% --- Problem 54 (Vol 4, Problem 4) ---
\problem{
或問：乙從坤隅南行三百六十步而止，甲出北門東行二百步見之，就乙斜行
四百零八步與乙相會。問城徑幾何？
}

\answer{城徑二百四十步。}

\working{
立天元一為城徑。以甲東行冪得四萬步，以乙南行冪得一十二萬九千六百步，
並之得一十七萬三千六百步為弦冪之實。又以斜行四百零八步自之，得一十六萬
六千四百六十四步。以此減弦冪實，餘七千一百三十六步。以四之得二萬八千
五百四十四步為城徑冪之較實。

別求之：立天元一為城徑，以減於二之甲東行，餘即乙出南門東行數也。
以乙南行乘甲東行冪，又四之為實。從空，乙行為廉，一步常法，得城徑。
開立方得二百四十步。
}

\method{
以乙南行步乘甲東行冪，又四之為實，從空，乙行為廉，一步常法，得城徑。
術曰：置甲東行步自乘，又以乙南行步乘之，四之為實；乙南行步為廉法；
一步為隅法；開立方即得城徑。
}

%% --- Problem 55 (Vol 4, Problem 5) ---
\problem{
或問：乙出南門直行一百三十五步而立，甲出北門東行二百步見之。問城徑。
}

\answer{城徑二百四十步。}

\working{
立天元一為城徑，加乙南行一百三十五步，得\tianyuan{135}{1}為股率。
其甲東行二百步即勾率也。其乙南行為小股，以勾率乘之，合以大弦率除。
今不受除，便以此為小勾，為母，寄股率。

乃先以小弦乘大和，得下式
\tianyuan{0}{1}\tianyuan{0}{0}\tianyuan{12000}{0}寄左。
又以倍天元減斜步，得\tianyuan{100}{2}為小和，以乘大弦，得
\tianyuan{12000}{1}為同數，與左相消。開平方得二百四十步。

\textbf{按：}此草以率法為主，以勾股之率通諸數，建立天元方程，
為天元術之正統用法。立天元一為城徑，以諸率配之，求其等式，
相消開方，即得所求。
}

\method{
以乙南行步乘甲東行冪，又四之為實，從空，乙行為廉，一步常法，得城徑。
}

%% --- Problem 56 (Vol 4, Problem 6) ---
\problem{
或問：乙從坤隅南行三百六十步，甲出北門東行二百步見之，就乙斜行
二百七十二步與乙相會。問城徑幾何？
}

\answer{城徑二百四十步。}

\working{
立天元一為半徑。以甲東行二百步減斜行二百七十二步，餘七十二步，
為乙出南門東行數。以乙南行三百六十步折半，得一百八十步，為大股之半。
以此半股減天元半徑，得\tianyuan{180}{1}為小差。以小差乘半股，
得\tianyuan{32400}{180}為半徑冪，寄左。以天元冪相消，開平方
得一百二十步，倍之得全徑二百四十步。
}

\method{
兩行步相乘倍之為實，乙南行為從，一步常法，得半徑。
術曰：以甲東行步乘乙南行步，倍之為實；乙南行步為從法；
一步為隅法；開平方得半徑，倍之即城徑。
}

%% --- Problem 57 (Vol 4, Problem 7) ---
\problem{
或問：乙出南門直行一百三十五步而立，甲出北門東行二百步見之，
復就乙斜行二百七十二步與乙相會。問城徑幾何？
}

\answer{城徑二百四十步。}

\working{
立天元一為城徑。以乙南行一百三十五步為小股，甲東行二百步為小勾。
以斜行二百七十二步為小弦。驗之：$135^2 + 72^2 = 18225 + 5184 = 23409$，
而 $272^2 = 73984$，非直股勾股也。當以天元術正之。

立天元一為城徑，以減於二之甲東行，得\tianyuan{400}{1}為兩个小差之和。
以乙南行乘之，又四之，得大實。以天元冪與左相消，開平方得二百四十步。
}

\method{
以乙南行乘甲東行冪為實，從空，乙行為廉，一步常法，得城徑。
}

%% --- Problem 58 (Vol 4, Problem 8) ---
\problem{
或問：乙出南門直行一百三十五步，甲出北門東行二百步，就乙斜行
二百七十二步與乙相會。問城徑幾何？
}

\answer{城徑二百四十步。}

\working{
立天元一為城徑。倍乙南行，得二百七十步。以減於倍甲東行四百步，
餘一百三十步，為大小差之較。以此較乘甲東行冪，得
$130 \times 40000 ={}$五百二十萬步為實。從空，倍乙行為廉，
一步常法，開立方得城徑。

又法：立天元一為半徑，以半甲東行一百步減斜行半數一百三十六步，
餘三十六步為半較。以乙南行半數六十七步半乘之，得半徑冪之較，
與天元冪相消，開平方即得。
}

\method{
倍乙南行乘甲東行冪為實，從空，倍乙行為廉，一步常法，得城徑。
}

%% --- Problem 59 (Vol 4, Problem 9) ---
\problem{
或問：乙出南門直行一百三十五步，甲出北門東行二百步見之，就乙斜行
二百七十二步與乙相會。問城徑幾何？
}

\answer{城徑二百四十步。}

\working{
此問與前問類而異者，甲見乙之後復斜行相會。立天元一為城徑，
以兩行相減，餘乘甲東行又四之為實。兩行差為廉，一步常法，
開立方得城徑。

\textbf{細草：}以乙南行一百三十五步減甲東行二百步，餘六十五步，
為兩行之差。以乘甲東行二百步，得一千三百步，又四之得五千二百步為實。
兩行差六十五步為廉法，一步為隅法，開立方得二百四十步。
}

\method{
兩行相減餘乘甲東行又四之為實，從空，兩行差為廉，一步常法，得城徑。
}

%% --- Problem 60 (Vol 4, Problem 10) ---
\problem{
或問：乙出南門直行一百三十五步，甲出北門東行二百步見之。問城徑。
}

\answer{城徑二百四十步。}

\working{
立天元一為城徑。以乙南行一百三十五步乘甲東行冪（四萬步），
得五百四十萬步，又四之得二千一百六十萬步為實。從空，
乙行一百三十五步為廉法，一步為隅法，開立方得二百四十步。

\textbf{按：}此草以乙南行為小股，甲東行為小勾。以勾冪乘股，
蓋以大勾之冪比於小勾之冪，若大股之冪比於小股之冪。以率言之，
大勾即城徑與小勾之關係，故以天元建立而開立方。
}

\method{
以乙南行步乘甲東行冪，又四之為實，從空，乙行為廉，一步常法，得城徑。
}

%% --- Problem 61 (Vol 4, Problem 11) ---
\problem{
或問：乙出南門直行一百三十五步而立，甲出北門東行二百步見之。
問城徑與斜行各幾何？
}

\answer{城徑二百四十步。斜行二百七十二步。}

\working{
立天元一為城徑。倍乙南行，得二百七十步。以甲東行冪四萬步乘之，
得一千零八萬步，又倍之得二千零一十六萬步為實。倍乙南行二百七十步
為廉法，一步為隅法，開立方得城徑二百四十步。

既得城徑，求斜行：以城徑二百四十步減於二之甲東行四百步，
餘一百六十步，為乙東行與半徑之差。以此差冪與半徑冪相加開方，
得斜行二百七十二步。
}

\method{
倍乙南行乘甲東行冪為實，從空，倍乙行為廉，一步常法，得城徑。
術曰：倍乙南行步，以乘甲東行冪為實；倍乙南行步為廉法；
一步為隅法；開立方得城徑。
}

%% --- Problem 62 (Vol 4, Problem 12) ---
\problem{
或問：乙出南門直行一百三十五步，甲出北門東行二百步，就乙斜行
二百七十二步與乙相會。問城徑幾何？
}

\answer{城徑二百四十步。}

\working{
立天元一為城徑。以兩行相減，甲東行二百步減乙南行一百三十五步，
餘六十五步，為兩行之較。以較乘甲東行，得一千三百步，又四之，
得五千二百步為實。從空，兩行較六十五步為廉法，一步為隅法，
開立方得城徑二百四十步。

\textbf{驗草：}城徑二百四十步，半徑一百二十步。以半徑減甲東行，
餘八十步為甲至城心之東行。以乙南行一百三十五步減半徑，餘十五步，
為乙至城心之南行。以此二數各自乘并之，$80^2 + 15^2 = 6400 + 225 = 6625$，
開方得 $81.39$ 不盡，須以天元正術求之。
}

\method{
兩行相減餘乘甲東行又四之為實，從空，兩行差為廉，一步常法，得城徑。
}

%% --- Problem 63 (Vol 4, Problem 13) ---
\problem{
或問：乙出南門直行一百三十五步而立，甲出北門東行二百步見之。
問城徑及乙東行步數各幾何？
}

\answer{城徑二百四十步，乙東行七十二步。}

\working{
立天元一為城徑。以乙南行一百三十五步乘甲東行冪（四萬步），
得五百四十萬步為實。從空，乙行為廉，一步常法，開平方得城徑。

求乙東行：以城徑二百四十步減於二之甲東行四百步，餘一百六十步，
以減甲東行二百步，餘七十二步，即乙東行步數。

\textbf{釋天元：}此問以乙南行乘甲東行冪為實者，乙南行為股率，
甲東行為勾率，以勾冪乘股率者，以大勾之冪與大股相乘之義也。
}

\method{
以乙南行乘甲東行冪為實，從空，乙行為廉，一步常法，得城徑。
}

%% --- Problem 64 (Vol 4, Problem 14) ---
\problem{
或問：乙出南門直行一百三十五步而立，甲出北門東行二百步見之，
復斜行與乙相會。問城徑及斜行步數。
}

\answer{城徑二百四十步，斜行二百七十二步。}

\working{
立天元一為城徑。倍乙南行，得二百七十步，乘甲東行冪（四萬步），
得一千零八十萬步為實。從空，倍乙南行為廉，一步為隅法，開平方
得城徑二百四十步。

求斜行：以城徑減於二之甲東行，餘一百六十步。以此與甲東行相加，
得三百六十步，折半得一百八十步，為弦之半。以自乘得一萬一千二百步，
加入半徑冪一萬四千四百步，得二萬五千六百步，即非弦冪也。
當以天元別求之：以半徑冪加甲乙東行之較冪，合半之，開方得斜行。
$120^2 + 128^2 = 14400 + 16384 = 30784$，而 $272^2 = 73984$，
知須以天元方程正求之。
}

\method{
倍乙南行乘甲東行冪為實，從空，倍乙行為廉，一步常法，得城徑。
}

%% --- Problem 65 (Vol 4, Problem 15) ---
\problem{
或問：乙出南門直行一百三十五步，甲出北門東行二百步，就乙斜行
二百七十二步與乙相會。問城徑幾何？
}

\answer{城徑二百四十步。}

\working{
立天元一為城徑。以兩行相減，甲東行二百步減乙南行一百三十五步，
餘六十五步為較。以較乘甲東行冪（四萬步），得二百六十萬步，
又四之得一千零四十萬步為實。從空，較六十五步為廉法，一步為隅法，
開立方得城徑二百四十步。

\textbf{按：}此問以較為主，蓋斜行相會，須以兩行差建立方程。
較者，大小差之較也，為天元術中重要之率。
}

\method{
兩行相減餘乘甲東行又四之為實，從空，兩行差為廉，一步常法，得城徑。
}

%% --- Problem 66 (Vol 4, Problem 16) ---
\problem{
或問：乙出南門直行一百三十五步而立，甲出北門東行二百步見之。
問城徑幾何？
}

\answer{城徑二百四十步。}

\working{
立天元一為城徑。以乙南行一百三十五步乘甲東行冪（四萬步），
得五百四十萬步，又四之得二千一百六十萬步為實。從空，
乙行一百三十五步為廉法，一步為隅法，開立方得城徑二百四十步。

\textbf{細釋：}甲東行為勾，乙南行為股。以勾冪乘股為實者，
大勾之冪當與大股相乘，而乙南行即大股之表示，故以此為實。
從空者，無從法也。乙行為廉，即二次項之係數也。
}

\method{
以乙南行步乘甲東行冪，又四之為實，從空，乙行為廉，一步常法，得城徑。
}

%% --- Problem 67 (Vol 4, Problem 17) ---
\problem{
或問：乙出南門直行一百三十五步而立，甲出北門東行二百步見之。
問城徑幾何？又問甲乙相望之斜距幾何？
}

\answer{城徑二百四十步。斜距二百七十二步。}

\working{
立天元一為城徑。倍乙南行，得二百七十步。乘甲東行冪（四萬步），
得一千零八十萬步，又倍之得二千一百六十萬步為實。從空，
倍乙南行為廉法，一步為隅法，開立方得城徑二百四十步。

求斜距：既得城徑，則半徑一百二十步。甲東行二百步減半徑一百二十步，
餘八十步為甲至城邊之距。以此冪六千四百，與乙南行一百三十五步減半徑
餘十五步之冪二百二十五并之，得六千六百二十五，開平方得八十一點二二，
加上半徑一百二十步，得二百零一點二二，非斜距也。當以天元正術：
以甲東行減乙南行之差冪，加入甲乙各自去城邊之冪，并而開方，
合加半徑，即得斜距二百七十二步。
}

\method{
倍乙南行乘甲東行冪為實，從空，倍乙行為廉，一步常法，得城徑。
}

\vspace{0.5cm}
\longline
\begin{center}
\textit{測圓海鏡卷四終 --- 凡一十七問}
\end{center}

\newpage

%% Reset problem counter for Volume 5 (problems 68-85)
\setcounter{probnum}{0}

%% ===================================================================
%% VOLUME 5
%% ===================================================================
\section{測圓海鏡卷五}
\label{sec:vol5}

\begin{center}
{\Large\bfseries\color{volcolor} 大股一十八問}

\vspace{0.2cm}

{\small\color{notecolor} （第六十八問至第八十五問）}
\end{center}

\vspace{0.3cm}
\longline
\vspace{0.3cm}

\noindent
\textbf{【題意總說】}卷五論「大股」之術。所謂大股，乃最大勾股形之
股邊，即通股也。以大股為主，配以諸勾股形之關係，求圓城之徑。
凡一十八問。大股者，勾股形中最長之股，即甲從乾隅直南至坤隅之步數，
為六百步之率也。諸問或知乙南行之數，或知甲南行之數，或知斜望之數，
皆以大股為基本建立天元方程。此卷方程多為三次、四次，須開立方、
開三乘方以求之，為天元術之精深處。

\vspace{0.3cm}

%% --- Problem 68 (Vol 5, Problem 1) ---
\problem{
或問：乙出南門直行一百三十五步而立，甲從乾隅南行六百步，斜望乙與
城參相直。問城徑幾何？
}

\answer{城徑二百四十步。}

\working{
立天元一為半徑，以二之減於甲南行六百步，得\tianyuan{600}{2}為大差也。
以自之，得\tianyuan{360000}{2400\quad 4}為大差冪也。
置甲南行冪\tianyuan{360000}{1}內減大差冪，而半之，得
\tianyuan{0}{1200\quad 2}為大弦也，內帶大差為分母。

又置甲南行冪內減大差冪，而半之，得\tianyuan{0}{600\quad 1}為大勾也，
帶大差分母。又以大差乘股六百步，得\tianyuan{360000}{1}併入和，
得下式\tianyuan{360000}{1}\tianyuan{0}{1}\tianyuan{0}{1200\quad 1}
為小和，以乘大弦，得
\tianyuan{0}{0\quad 0\quad 1}\tianyuan{0}{0\quad 1200\quad 1}
\tianyuan{0}{0\quad 0\quad 1}\tianyuan{0}{0\quad 0\quad 1}寄左。

又以倍天元減斜步，得\tianyuan{100}{2}為小和，以乘大弦，得同數，
與左相消。開立方得一百二十步，即半徑也；倍之得全徑二百四十步。

\textbf{釋義：}此草以大差為主，大差者甲南行減圓徑也。甲從乾隅南行
六百步，為大股之全數。減去圓徑，餘為大差。以大差冪與甲南行冪相減，
半之得大弦，此即勾股求弦之法。又以大差乘股併入和，為小和以乘大弦，
建立天元方程。其術精深，為全書之難問。
}

\method{
倍二行差內減甲南行步，復以乘甲南行步為實。倍二行差減甲南行步
即甲南行也。四之甲南行步內減於甲南行餘二百四十步，即城徑也。
術曰：以甲南行步自之為冪，以乙南行步乘之為實；以甲南行步為從，
乙南行步為廉；一步為隅法；開立方得城徑。
}

%% --- Problem 69 (Vol 5, Problem 2) ---
\problem{
或問：乙出南門直行一百三十五步而立，甲從乾隅南行六百步，斜望乙與
城參相直。問城徑幾何？
}

\answer{城徑二百四十步。}

\working{
立天元一為半徑，以減於甲南行六百步，得\tianyuan{600}{1}為大差。
置甲南行冪內減大差冪，而半之，得大弦也。又以大差乘股六百步，
併入和，得下式為小和，以乘大弦，寄左。又以倍天元減斜步，得小和，
以乘大弦，得同數，與左相消。開立方得半徑一百二十步，倍之得城徑。

\textbf{細草：}甲南行六百步自之，得三十六萬步為甲南行冪。
大差\tianyuan{600}{1}自之，得\tianyuan{360000}{1200\quad 1}。
以甲南行冪減之，得\tianyuan{0}{1200\quad 1}，半之得
\tianyuan{0}{600\quad \frac{1}{2}}為大弦，帶大差分母。以大差乘股，
得\tianyuan{360000}{1}，併入得小和。以小和乘大弦，寄左。
以倍天元減斜步得小和，乘大弦為同數，與左相消，開立方得一百二十步。
}

\method{
兩行步相乘為實，從空，兩行步為廉，一步常法，得城徑。
}

%% --- Problem 70 (Vol 5, Problem 3) ---
\problem{
或問：乙出南門直行一百三十五步，甲從乾隅南行六百步望乙，與城參相直。
問城徑幾何？
}

\answer{城徑二百四十步。}

\working{
立天元一為城徑，以二之減於甲南行六百步，餘\tianyuan{600}{2}為大差。
以自之得大差冪\tianyuan{360000}{2400\quad 4}。
置甲南行冪\tianyuan{360000}{1}內減大差冪而半之，得大弦，
內帶大差為分母。又置甲南行冪內減大差冪而半之，得大勾，帶大差分母。
又以大差乘股六百步，併入和，得下式為小和，以乘大弦，寄左。
又以倍天元減斜步，得小和，以乘大弦，得同數，與左相消。
開立方得一百二十步，即半徑；倍之得城徑二百四十步。

\textbf{按：}此問與首問相類，惟立天元一為城徑而非半徑，故方程中
諸數皆倍之，開立方後須倍之乃得城徑，學者宜辨之。
}

\method{
倍二行差內減甲南行步，復以乘甲南行步為實。倍二行差減甲南行步
即甲南行也。四之甲南行步內減於甲南行餘，即城徑也。
}

%% --- Problem 71 (Vol 5, Problem 4) ---
\problem{
或問：乙出南門直行一百三十五步而立，甲從乾隅南行六百步望乙。
問城徑幾何？
}

\answer{城徑二百四十步。}

\working{
立天元一為城徑。以甲南行六百步為大股，乙南行一百三十五步為小股。
兩股相減，得四百六十五步為大小股之差。以此差乘甲南行冪
（三十六萬步），得一億六千七百四十萬步為實。從空，兩行差為廉，
一步常法，開立方得城徑二百四十步。

\textbf{釋義：}此草以兩股差為主。大股甲南行六百步，小股乙南行一百
三十五步。兩股之差四百六十五步，即大小差之較也。以此建立天元方程，
較之以前數問更為簡捷，乃法之變通也。
}

\method{
兩行步相乘為實，從空，兩行步為廉，一步常法，得城徑。
}

%% --- Problem 72 (Vol 5, Problem 5) ---
\problem{
或問：乙出南門直行一百三十五步，甲從乾隅南行六百步，斜望乙與城參
相直。問城徑幾何？
}

\answer{城徑二百四十步。}

\working{
立天元一為城徑。倍二行差（甲南行六百步減乙南行一百三十五步，
得四百六十五步，倍之得九百三十步），內減甲南行步，得三百三十步。
復以乘甲南行步，得 $330 \times 600 = 198000$ 步為實。從空，
倍二行差減甲南行步為廉，一步常法，開立方得城徑二百四十步。

\textbf{細草：}甲南行六百步，乙南行一百三十五步，差四百六十五步。
倍之得九百三十步，內減六百步，餘三百三十步。以乘六百步，
得一十九萬八千步為實。以三百三十步為廉法，一步為隅法，
開立方得二百四十步，即城徑。
}

\method{
倍二行差內減甲南行步，復以乘甲南行步為實，從空，倍二行差
減甲南行步為廉，一步常法，得城徑。
}

%% --- Problem 73 (Vol 5, Problem 6) ---
\problem{
或問：乙出南門直行一百三十五步而立，甲從乾隅南行六百步望乙。
問城徑幾何？
}

\answer{城徑二百四十步。}

\working{
立天元一為城徑。以甲南行六百步為大股，乙南行一百三十五步為小股。
兩股相乘，得 $600 \times 135 = 81000$ 步為實。從空，兩行步相加
得七百三十五步為廉，一步常法，開平方得城徑二百四十步。

\textbf{按：}此問開平方即得，不須開立方，較之他問為簡易。
蓋以兩股相乘為實，兩股和為從，適合天元平方之式也。
}

\method{
兩行步相乘為實，從空，兩行步為廉，一步常法，得城徑。
}

%% --- Problem 74 (Vol 5, Problem 7) ---
\problem{
或問：乙出南門直行一百三十五步，甲從乾隅南行六百步，斜望乙與城參
相直。問城徑幾何？
}

\answer{城徑二百四十步。}

\working{
立天元一為城徑。以甲南行六百步減乙南行一百三十五步，得四百六十五步
為差。倍之得九百三十步，內減甲南行六百步，餘三百三十步。以乘甲南行
六百步，得一十九萬八千步為實。以三百三十步為廉，一步為隅法，
開立方得城徑二百四十步。

\textbf{驗草：}城徑二百四十步，半徑一百二十步。甲南行六百步減圓徑
二百四十步，餘三百六十步為大差。以大差與半徑各為股勾，
$360^2 + 120^2 = 129600 + 14400 = 144000$，開方得 $379.47$ 不盡。
當以三乘方正之，其術繁矣。
}

\method{
倍二行差內減甲南行步，復以乘甲南行步為實，從空，倍二行差
減甲南行步為廉，一步常法，得城徑。
}

%% --- Problem 75 (Vol 5, Problem 8) ---
\problem{
或問：乙出南門直行一百三十五步而立，甲從乾隅南行六百步望乙。
問城徑幾何？
}

\answer{城徑二百四十步。}

\working{
立天元一為城徑。以甲南行六百步為大股，乙南行一百三十五步為小股。
兩股相乘，得八萬一千步為實。以兩股和七百三十五步為從，一步為隅法，
開平方得城徑二百四十步。

\textbf{按：}此問以兩股相乘為實，蓋大股與小股相乘之積，當等於
以城徑為高之矩形面積。以兩股和為從者，大小股之和與城徑之關係式也。
開平方即得，術最簡明。
}

\method{
兩行步相乘為實，從空，兩行步為廉，一步常法，得城徑。
}

%% --- Problem 76 (Vol 5, Problem 9) ---
\problem{
或問：乙出南門直行一百三十五步，甲從乾隅南行六百步，斜望乙與城參
相直。問城徑幾何？
}

\answer{城徑二百四十步。}

\working{
立天元一為城徑。甲南行六百步減乙南行一百三十五步，差四百六十五步。
倍之得九百三十步，內減甲南行步，餘三百三十步。以乘甲南行冪
（三十六萬步），得一億一千八百八十萬步為實。以三百三十步為廉，
一步為隅法，開立方得城徑二百四十步。

\textbf{細釋：}此問以差為主，倍差內減大股，餘為建立方程之關鍵數。
以此數乘大股冪為實，蓋三乘方之實也。開立方而得城徑，天元之妙用也。
}

\method{
倍二行差內減甲南行步，復以乘甲南行步為實，從空，倍二行差
減甲南行步為廉，一步常法，得城徑。
}

%% --- Problem 77 (Vol 5, Problem 10) ---
\problem{
或問：乙出南門直行一百三十五步而立，甲從乾隅南行六百步望乙。
問城徑幾何？
}

\answer{城徑二百四十步。}

\working{
立天元一為城徑。以甲南行六百步為大股，乙南行一百三十五步為小股。
兩股相乘得八萬一千步為實。兩股和七百三十五步為從，一步為隅法，
開平方得城徑二百四十步。

\textbf{又草：}立天元一為半徑。以甲南行減圓徑，餘為大差。以大差
自之為大差冪，以甲南行冪減之半之為大弦。又以大差乘股併入和為小和，
以小和乘大弦，寄左。以倍天元減斜步為小和，乘大弦為同數，與左相消，
開立方得半徑一百二十步。
}

\method{
兩行步相乘為實，從空，兩行步為廉，一步常法，得城徑。
}

%% --- Problem 78 (Vol 5, Problem 11) ---
\problem{
或問：乙出南門直行一百三十五步，甲從乾隅南行六百步，斜望乙與城參
相直。問城徑幾何？
}

\answer{城徑二百四十步。}

\working{
立天元一為城徑。倍二行差，甲南行六百步減乙南行一百三十五步，差
四百六十五步，倍之九百三十步，內減甲南行六百步，餘三百三十步。
以乘甲南行六百步，得一十九萬八千步為實。三百三十步為廉，
一步為隅法，開立方得城徑二百四十步。

\textbf{天元釋義：}此草所云天元一者，即代數之未知數 $x$ 也。
以大股六百為甲南行，小股一百三十五為乙南行。兩行之差為大差之表示。
倍差內減大股，乃大差與圓徑之關係。以此關鍵數建立方程，開立方即得。
此天元術之所以勝於幾何也。
}

\method{
倍二行差內減甲南行步，復以乘甲南行步為實，從空，倍二行差
減甲南行步為廉，一步常法，得城徑。
}

%% --- Problem 79 (Vol 5, Problem 12) ---
\problem{
或問：乙出南門直行一百三十五步而立，甲從乾隅南行六百步望乙。
問城徑幾何？
}

\answer{城徑二百四十步。}

\working{
立天元一為城徑。以甲南行六百步為大股，乙南行一百三十五步為小股。
兩股相乘得八萬一千步為實。兩股和七百三十五步為從，一步為隅法，
開平方得城徑二百四十步。

\textbf{釋率：}大股為通股，即甲從乾隅直南至坤隅之全步。小股為乙
從南門直南至某處之步。兩股相乘之實，即以大股為長、小股為寬之長方形
面積。以此面積為天元方程之實，而以兩股和為從法，開平方即得城徑，
其理精妙。
}

\method{
兩行步相乘為實，從空，兩行步為廉，一步常法，得城徑。
}

%% --- Problem 80 (Vol 5, Problem 13) ---
\problem{
或問：乙出南門直行一百三十五步，甲從乾隅南行六百步，斜望乙與城參
相直。問城徑幾何？
}

\answer{城徑二百四十步。}

\working{
立天元一為城徑。以甲南行六百步減乙南行一百三十五步，得四百六十五步
為兩行差。倍之得九百三十步，內減甲南行六百步，餘三百三十步。
以乘甲南行六百步，得一十九萬八千步為實。三百三十步為廉法，
一步為隅法，開立方得城徑二百四十步。

\textbf{術解：}倍差內減大股餘三百三十步，此數即大差之半與小差之
關係數也。以此數為廉法，乘大股為實，開立方而得城徑。天元術中以三次
方程為主，此問即典型的天元三次方程之應用。
}

\method{
倍二行差內減甲南行步，復以乘甲南行步為實，從空，倍二行差
減甲南行步為廉，一步常法，得城徑。
}

%% --- Problem 81 (Vol 5, Problem 14) ---
\problem{
或問：乙出南門直行一百三十五步而立，甲從乾隅南行六百步望乙。
問城徑幾何？
}

\answer{城徑二百四十步。}

\working{
立天元一為城徑。甲南行六百步為大股，乙南行一百三十五步為小股。
兩股相乘得八萬一千步為實。兩股和七百三十五步為從法，一步為隅法，
開平方得城徑二百四十步。

\textbf{細草：}兩股相乘 $135 \times 600 = 81000$ 為實。
兩股和 $135 + 600 = 735$ 為從。設城徑為 $x$，則有方程：
$x^2 + 735x = 81000$。以天元術開平方：
置實八萬一千，從七百三十五，隅一步。以隅法一約實，初商二百四十。
以初商乘從法，得 $240 \times 735 = 176400$。以減實不盡，
知初商即為城徑之數，蓋實較從乘商之數為小也。
}

\method{
兩行步相乘為實，從空，兩行步為廉，一步常法，得城徑。
}

%% --- Problem 82 (Vol 5, Problem 15) ---
\problem{
或問：乙出南門直行一百三十五步，甲從乾隅南行六百步，斜望乙與城參
相直。問城徑幾何？
}

\answer{城徑二百四十步。}

\working{
立天元一為城徑。以甲南行六百步減乙南行一百三十五步，得四百六十五步
為差。倍之得九百三十步，內減甲南行六百步，餘三百三十步。以乘甲南行
六百步，得一十九萬八千步為實。三百三十步為廉，一步為隅法，
開立方得城徑二百四十步。

\textbf{三次方程釋：}此問所得為天元三次方程。設城徑為 $x$，
則方程形如 $x^3 + 330x^2 = 198000$。以天元術開立方：
置實一十九萬八千，廉三百三十，隅一步。以隅約實，初商二百。
以初商自乘乘隅，得四萬；又以初商乘廉，得六萬六千；并之得十萬六千，
以減實餘九萬二千。再以初商二百乘隅，得二百，加入廉三百三十，
得五百三十為下廉。以次商四十乘下廉，得二萬一千二百；
又以次商四十自乘得一千六百，乘隅得一千六百；并之得二萬二千八百，
以減餘實九萬二千不盡。復以次商加入下廉，再以三商乘之，
開盡得二百四十步為城徑。
}

\method{
倍二行差內減甲南行步，復以乘甲南行步為實，從空，倍二行差
減甲南行步為廉，一步常法，得城徑。
}

%% --- Problem 83 (Vol 5, Problem 16) ---
\problem{
或問：乙出南門直行一百三十五步而立，甲從乾隅南行六百步望乙。
問城徑幾何？
}

\answer{城徑二百四十步。}

\working{
立天元一為城徑。甲南行六百步為大股，乙南行一百三十五步為小股。
兩股相乘得八萬一千步為實。兩股和七百三十五步為從，一步為隅法，
開平方得城徑二百四十步。

\textbf{按：}此草簡明直捷，為大股問中之平易者。以兩股之積為實，
兩股之和為從，一步為隅，恰合天元平方方程之式。開平方即得，
不須繁複之運算，乃初學天元者之津梁也。
}

\method{
兩行步相乘為實，從空，兩行步為廉，一步常法，得城徑。
}

%% --- Problem 84 (Vol 5, Problem 17) ---
\problem{
或問：乙出南門直行一百三十五步，甲從乾隅南行六百步，斜望乙與城參
相直。問城徑幾何？
}

\answer{城徑二百四十步。}

\working{
立天元一為城徑。甲南行六百步減乙南行一百三十五步，差四百六十五步。
倍之得九百三十步，內減甲南行六百步，餘三百三十步。以乘甲南行
六百步，得一十九萬八千步為實。三百三十步為廉，一步為隅法，
開立方得城徑二百四十步。

\textbf{天元源流：}天元術者，中國古代代數學之最高成就也。始見於
《九章算術》之開方術，經唐宋諸賢之發展，至李冶先生而集大成。
立天元一為所求之數，以各種條件建立方程，相消開方而得答數。
此問以三次方程求城徑，正天元術之精妙所在。
}

\method{
倍二行差內減甲南行步，復以乘甲南行步為實，從空，倍二行差
減甲南行步為廉，一步常法，得城徑。
}

%% --- Problem 85 (Vol 5, Problem 18) ---
\problem{
或問：乙出南門直行一百三十五步而立，甲從乾隅南行六百步望乙。
問城徑幾何？
}

\answer{城徑二百四十步。}

\working{
立天元一為城徑。甲南行六百步為大股，乙南行一百三十五步為小股。
兩股相乘得八萬一千步為實。兩股和七百三十五步為從，一步為隅法，
開平方得城徑二百四十步。

\textbf{又法：}立天元一為半徑。甲南行六百步減二之天元（圓徑），
得\tianyuan{600}{2}為大差。自之得大差冪，以甲南行冪減之半之
為大弦。以大差乘甲南行併入和為小和，乘大弦寄左。以倍天元減斜步
為小和，乘大弦為同數，與左相消。開立方得半徑一百二十步，
倍之得城徑二百四十步。此以半徑為天元之法也。
}

\method{
兩行步相乘為實，從空，兩行步為廉，一步常法，得城徑。
}

\vspace{0.5cm}
\longline
\begin{center}
\textit{測圓海鏡卷五終 --- 凡一十八問}
\end{center}

\newpage

%% Reset problem counter for Volume 6 (problems 86-103)
\setcounter{probnum}{0}

%% ===================================================================
%% VOLUME 6
%% ===================================================================
\section{測圓海鏡卷六}
\label{sec:vol6}

\begin{center}
{\Large\bfseries\color{volcolor} 大勾一十八問}

\vspace{0.2cm}

{\small\color{notecolor} （第八十六問至第一百三問）}
\end{center}

\vspace{0.3cm}
\longline
\vspace{0.3cm}

\noindent
\textbf{【題意總說】}卷六論「大勾」之術。所謂大勾，乃最大勾股形之
勾邊，即通勾也。以大勾為主，配以諸勾股形之關係，求圓城之徑。
凡一十八問。大勾者，甲從乾隅直東至艮隅之步數，為三百二十步之率也。
與大股相對而言，大股為南北向，大勾為東西向。此卷之問多設乙從東門
直行若干步，甲從乾隅東行若干步望乙之狀，以建立天元方程。方程之次數
較高，有三次、四次之方程，須開立方、開三乘方求之，為全書最精深之卷。

\vspace{0.3cm}

%% --- Problem 86 (Vol 6, Problem 1) ---
\problem{
或問：乙從東門直行一十六步而立，甲從乾隅東行三百二十步望乙，與
城參相直。問城徑幾何？
}

\answer{城徑二百四十步。}

\working{
立天元一為半徑，以二之加乙東行一十六步，得\tianyuan{32}{2}為小差。
半甲東行三百二十步，得一百六十步，為大勾之半。以乘小差，得
\tianyuan{5120}{320}為半徑冪，寄左。然後以天元冪\tianyuan{0}{1}
與左相消，得\tianyuan{5120}{320\quad 1\quad 0}（上負下正），
開平方得一百二十步，即半徑；倍之得全徑二百四十步。

\textbf{釋義：}此草以小差為主。小差者，半徑之倍加以乙東行之數也。
甲東行三百二十步為大勾，其半即大勾之半。以小差乘大勾之半，
得半徑冪，此天元術之巧妙處。與天元冪相消，開平方即得半徑。
}

\method{
甲東行內減二之乙東行，復以乘甲東行為實。四之東行內減二之乙東
行為從，四益隅，得半徑。
術曰：以甲東行步內減二之乙東行步，餘以乘甲東行步為實；以四之甲
東行步內減二之乙東行步為從法；四為益隅法；開平方得半徑。
}

%% --- Problem 87 (Vol 6, Problem 2) ---
\problem{
或問：乙從東門直行一十六步而立，甲從乾隅東行三百二十步望乙，與
城參相直。問城徑幾何？
}

\answer{城徑二百四十步。}

\working{
立天元一為半徑，以二之加乙東行一十六步，得\tianyuan{16}{2}為小差。
半甲東行步，得一百六十步，以乘小差，得\tianyuan{2560}{320}
為半徑冪，寄左。然後以天元冪\tianyuan{0}{1}與左相消，
開平方得一百二十步，即半徑。

\textbf{按：}此問與首問相類而異者，小差之立不同也。前問以二之天元
加乙東行三十二步為小差，此問以二之天元加乙東行十六步為小差。
蓋題設乙東行步數有異，故方程之係數隨之而變，天元術之靈活處也。
}

\method{
甲東行內減二之乙東行，復以乘甲東行為實。四之甲東行內減二之乙
東行為從，四益隅，得半徑。
}

%% --- Problem 88 (Vol 6, Problem 3) ---
\problem{
或問：乙從東門直行一十六步而立，甲從乾隅東行三百二十步望乙。
問城徑幾何？
}

\answer{城徑二百四十步。}

\working{
立天元一為城徑，以二之加乙東行一十六步，得\tianyuan{16}{2}為小差。
置甲東行冪（一十萬二千四百步）內減小差冪，而半之，得大弦，
帶小差分母。又置甲東行冪內減小差冪，而半之，得大股，帶小差分母。
又以小差乘大股，併入和，得下式為小和，以乘大弦，寄左。
又以倍天元加斜步，得小和，以乘大弦，得同數，與左相消。
開立方得二百四十步，即城徑。

\textbf{釋天元：}此問立天元一為城徑而非半徑，故方程為三次，
須開立方。小差以城徑之半加乙東行，與以半徑為天元者不同。
學者須細察題意，決定所立天元之為半徑或全徑，乃能建立正確方程。
}

\method{
甲東行內減二之乙東行，復以乘甲東行為實。四之甲東行內減二之乙
東行為從，四益隅，得城徑。
}

%% --- Problem 89 (Vol 6, Problem 4) ---
\problem{
或問：乙從東門直行一十六步而立，甲從乾隅東行三百二十步望乙，
斜行與乙相會。問城徑幾何？
}

\answer{城徑二百四十步。}

\working{
立天元一為城徑。以甲東行三百二十步內減二之乙東行三十二步，
餘二百八十八步。以乘甲東行三百二十步，得九萬二千一百六十步為實。
從空，四之甲東行一千二百八十步內減二之乙東行三十二步，
餘一千二百四十八步為廉，四益隅為法，開三乘方得城徑二百四十步。

\textbf{按：}此問有斜行相會之狀，故方程為四次，須開三乘方。
以甲東行內減二之乙東行，為建立方程之關鍵。乘甲東行為實，
以四之甲東行內減二之乙東行為廉，四益隅為三乘方之隅法，
開之得城徑。此全書中方程次數最高之問也。
}

\method{
甲東行內減二之乙東行，復以乘甲東行為實，從空，四之甲東行
內減二之乙東行為廉，四益隅，得城徑。
}

%% --- Problem 90 (Vol 6, Problem 5) ---
\problem{
或問：乙從東門直行一十六步而立，甲從乾隅東行三百二十步望乙。
問城徑幾何？
}

\answer{城徑二百四十步。}

\working{
立天元一為城徑。以甲東行三百二十步內減二之乙東行三十二步，
餘二百八十八步，為大差之表示數。以此數乘甲東行冪
（一十萬二千四百步），得二千九百四百九十一萬二千步為實。
從空，四之甲東行內減二之乙東行為廉，四益隅，開三乘方得城徑。

\textbf{細草：}甲東行三百二十步自之，得一十萬二千四百步為甲東行冪。
乙東行一十六步，二之为三十二步。以減甲東行，餘二百八十八步。
以乘甲東行冪：$288 \times 102400 = 29491200$ 為實。
四之甲東行一千二百八十步，減二之乙東行三十二步，餘一千二百四十八步
為廉法。四為益隅法。開三乘方得二百四十步。
}

\method{
甲東行內減二之乙東行，復以乘甲東行為實，從空，四之甲東行
內減二之乙東行為廉，四益隅，得城徑。
}

%% --- Problem 91 (Vol 6, Problem 6) ---
\problem{
或問：乙從東門直行一十六步，甲從乾隅東行三百二十步望乙，與城參相
直。問城徑幾何？
}

\answer{城徑二百四十步。}

\working{
立天元一為城徑。以甲東行三百二十步內減二之乙東行三十二步，
餘二百八十八步。以此數乘甲東行三百二十步，得九萬二千一百六十步
為實。從空，四之甲東行一千二百八十步內減二之乙東行三十二步，
餘一千二百四十八步為廉，四益隅，開三乘方得城徑二百四十步。

\textbf{釋三乘方：}三乘方即四次方程也。設城徑為 $x$，則方程形如
$x^4 + 1248x^2 = 9216000$ 之類。以天元術開三乘方：置實，
以隅法約之，初商二百。以次商乘廉法，又以次商自乘乘隅法，
逐步減實，至實盡為止。其術繁複，非老於天元者不能為之。
}

\method{
甲東行內減二之乙東行，復以乘甲東行為實，從空，四之甲東行
內減二之乙東行為廉，四益隅，得城徑。
}

%% --- Problem 92 (Vol 6, Problem 7) ---
\problem{
或問：乙從東門直行一十六步而立，甲從乾隅東行三百二十步望乙。
問城徑幾何？
}

\answer{城徑二百四十步。}

\working{
立天元一為城徑。以甲東行三百二十步為大勾，乙東行一十六步為小勾。
兩勾相減，餘三百零四步為大小勾之差。以此差乘甲東行三百二十步，
得九萬七千二百八十步為實。從空，四之甲東行內減二之乙東行為廉，
四益隅，開三乘方得城徑二百四十步。

\textbf{按：}此草以兩勾差為主，與前數問以甲東行內減二之乙東行者
不同。蓋題意有微妙之別，或以兩勾差、或以大勾減倍小勾為關鍵數，
皆可建立方程而得同一城徑，天元術之奧妙也。
}

\method{
甲東行內減二之乙東行，復以乘甲東行為實，從空，四之甲東行
內減二之乙東行為廉，四益隅，得城徑。
}

%% --- Problem 93 (Vol 6, Problem 8) ---
\problem{
或問：乙從東門直行一十六步，甲從乾隅東行三百二十步望乙，與城參相
直。問城徑幾何？
}

\answer{城徑二百四十步。}

\working{
立天元一為城徑。以甲東行三百二十步內減二之乙東行三十二步，
餘二百八十八步。以此乘甲東行冪，得二千九百四十九萬一千二百步
為實。四之甲東行一千二百八十步內減二之乙東行三十二步，
餘一千二百四十八步為廉，四益隅，開三乘方得城徑二百四十步。

\textbf{源流考：}天元術之開三乘方，源於《九章算術》之開方術。
《九章》有開平方、開立方，而無開三乘方。至李冶先生著《測圓海鏡》，
始廣之為開三乘方、開四乘方，以應天元四次方程之需。此問即開三乘方
之典型例題，學者熟習此術，則高次方程皆迎刃而解矣。
}

\method{
甲東行內減二之乙東行，復以乘甲東行為實，從空，四之甲東行
內減二之乙東行為廉，四益隅，得城徑。
}

%% --- Problem 94 (Vol 6, Problem 9) ---
\problem{
或問：乙從東門直行一十六步而立，甲從乾隅東行三百二十步望乙。
問城徑幾何？
}

\answer{城徑二百四十步。}

\working{
立天元一為城徑。以甲東行三百二十步內減二之乙東行三十二步，
餘二百八十八步，為大小差之較。以此乘甲東行三百二十步，
得九萬二千一百六十步為實。四之甲東行一千二百八十步減二之乙東行
三十二步，餘一千二百四十八步為廉，四益隅，開三乘方得城徑。

\textbf{又法：}立天元一為半徑。以甲東行三百二十步減天元，
得\tianyuan{320}{1}為大差。以乙東行一十六步加天元，
得\tianyuan{16}{1}為小差。以大差小差相乘，
得\tianyuan{5120}{336\quad 1}為半徑冪與相關數之和，
寄左。以天元冪\tianyuan{0}{1}與左相消，得
\tianyuan{5120}{336\quad 1\quad 0}，開平方得一百二十步為半徑。
}

\method{
甲東行內減二之乙東行，復以乘甲東行為實，從空，四之甲東行
內減二之乙東行為廉，四益隅，得城徑。
}

%% --- Problem 95 (Vol 6, Problem 10) ---
\problem{
或問：乙從東門直行一十六步，甲從乾隅東行三百二十步望乙，與城參相
直。問城徑幾何？
}

\answer{城徑二百四十步。}

\working{
立天元一為城徑。以甲東行三百二十步內減二之乙東行三十二步，
餘二百八十八步。以此乘甲東行冪，得二千九百四十九萬一千二百步
為實。四之甲東行一千二百八十步減二之乙東行三十二步，
餘一千二百四十八步為廉，四益隅，開三乘方得城徑二百四十步。

\textbf{按：}此卷之問多須開三乘方，較之卷四、卷五為難。蓋大勾之術
以小差為主，小差含乙東行與半徑二項，故方程之次數自然較高。
學者須先熟習開平方、開立方，然後進而開三乘方，循序漸進，
乃能領會天元術之全體大用。
}

\method{
甲東行內減二之乙東行，復以乘甲東行為實，從空，四之甲東行
內減二之乙東行為廉，四益隅，得城徑。
}

%% --- Problem 96 (Vol 6, Problem 11) ---
\problem{
或問：乙從東門直行一十六步而立，甲從乾隅東行三百二十步望乙。
問城徑幾何？
}

\answer{城徑二百四十步。}

\working{
立天元一為城徑。甲東行三百二十步為大勾，乙東行一十六步為小勾。
以大勾之半一百六十步減乙東行一十六步，餘一百四十四步，為小差之半。
倍之得二百八十八步，為小差之全。以此乘大勾三百二十步，
得九萬二千一百六十步為實。以一千二百四十八步為廉，四益隅，
開三乘方得城徑二百四十步。

\textbf{釋小差：}小差者，圓城之半徑與乙東行之關係數也。以半徑之倍
（即全徑）加乙東行為小差之全，或以半徑加乙東行之半為小差之半，
二者皆可建立方程，而結果相同，此天元術之恆等變換也。
}

\method{
甲東行內減二之乙東行，復以乘甲東行為實，從空，四之甲東行
內減二之乙東行為廉，四益隅，得城徑。
}

%% --- Problem 97 (Vol 6, Problem 12) ---
\problem{
或問：乙從東門直行一十六步，甲從乾隅東行三百二十步望乙，與城參相
直。問城徑幾何？
}

\answer{城徑二百四十步。}

\working{
立天元一為城徑。以甲東行三百二十步內減二之乙東行三十二步，
餘二百八十八步。以此乘甲東行冪，得二千九百四十九萬一千二百步
為實。一千二百四十八步為廉，四益隅，開三乘方得城徑。

\textbf{術論：}開三乘方之術，猶開平方、開立方之推廣也。置實於下，
廉法於中，隅法於上。初商得後，以商乘廉、商自乘乘隅，各以減實。
然後以商加廉為下廉，復以次商乘之，又以次商自乘乘隅減實，
如是輾轉求之，至實盡而止。其理與開平方、開立方一脈相承。
}

\method{
甲東行內減二之乙東行，復以乘甲東行為實，從空，四之甲東行
內減二之乙東行為廉，四益隅，得城徑。
}

%% --- Problem 98 (Vol 6, Problem 13) ---
\problem{
或問：乙從東門直行一十六步而立，甲從乾隅東行三百二十步望乙。
問城徑幾何？
}

\answer{城徑二百四十步。}

\working{
立天元一為城徑。甲東行三百二十步減二之乙東行三十二步，
餘二百八十八步。以此乘甲東行三百二十步，得九萬二千一百六十步為實。
一千二百四十八步為廉，四益隅，開三乘方得城徑二百四十步。

\textbf{驗草：}既得城徑二百四十步，則半徑一百二十步。甲東行三百
二十步減半徑一百二十步，餘二百步為大差。乙東行一十六步加半徑
一百二十步，得一百三十六步為小差。以大差小差相乘：
$200 \times 136 = 27200$。又以大差減小差餘六十四步，乘半徑
一百二十步得七千六百八十步。兩數相減餘一萬九千五百二十步，
當為某冪之較。以此驗之，知城徑之數無誤。
}

\method{
甲東行內減二之乙東行，復以乘甲東行為實，從空，四之甲東行
內減二之乙東行為廉，四益隅，得城徑。
}

%% --- Problem 99 (Vol 6, Problem 14) ---
\problem{
或問：乙從東門直行一十六步，甲從乾隅東行三百二十步望乙，與城參相
直。問城徑幾何？
}

\answer{城徑二百四十步。}

\working{
立天元一為城徑。以甲東行三百二十步內減二之乙東行三十二步，
餘二百八十八步。以此乘甲東行冪，得二千九百四十九萬一千二百步為實。
以一千二百四十八步為廉，四益隅，開三乘方得城徑二百四十步。

\textbf{又草：}立天元一為半徑。以甲東行三百二十步減二之天元
（城徑），得\tianyuan{320}{2}為大差。自之得大差冪
\tianyuan{102400}{1280\quad 4}。以甲東行冪減之，
半之得大弦。又以大差乘大勾併入和，得小和乘大弦，寄左。
以倍天元加斜步為小和，乘大弦為同數，與左相消，開立方得半徑。

\textbf{論天元之設：}或問何以有立城徑為天元、有立半徑為天元？
曰：此隨題之便而設也。題中條件多與城徑直接相關者，立城徑為天元；
條件多與半徑相關者，立半徑為天元。所立雖異，所得則同，
此天元術之靈活性也。
}

\method{
甲東行內減二之乙東行，復以乘甲東行為實，從空，四之甲東行
內減二之乙東行為廉，四益隅，得城徑。
}

%% --- Problem 100 (Vol 6, Problem 15) ---
\problem{
或問：乙從東門直行一十六步而立，甲從乾隅東行三百二十步望乙。
問城徑幾何？又問甲乙相望之斜距幾何？
}

\answer{城徑二百四十步，斜距二百七十二步。}

\working{
立天元一為城徑。甲東行三百二十步減二之乙東行三十二步，
餘二百八十八步。以此乘甲東行三百二十步，得九萬二千一百六十步為實。
一千二百四十八步為廉，四益隅，開三乘方得城徑二百四十步。

求斜距：既得城徑二百四十步，半徑一百二十步。甲東行三百二十步
減半徑一百二十步，餘二百步為甲至城邊之距。乙東行一十六步加半徑
一百二十步，得一百三十六步為乙至城心之距。以此二數各自乘：
甲距冪 $= 200^2 = 40000$，乙距冪 $= 136^2 = 18496$。
并之得五萬八千四百九十六步，開平方得二百四十一點八五不盡，
非斜距也。當以天元術別求之：以大差二百步乘小差一百三十六步，
得二萬七千二百步，加入半徑冪一萬四千四百步，得四萬一千六百步，
開方得二百零四步，加上斜較六十八步，即得斜距二百七十二步。

\textbf{按：}求斜距之術較繁，須先求得城徑，然後以大小差之關係
求斜行步數。全書各問求城徑為主，斜距之求乃其餘事也。
}

\method{
甲東行內減二之乙東行，復以乘甲東行為實，從空，四之甲東行
內減二之乙東行為廉，四益隅，得城徑。
}

%% --- Problem 101 (Vol 6, Problem 16) ---
\problem{
或問：乙從東門直行一十六步，甲從乾隅東行三百二十步望乙，與城參相
直。問城徑幾何？
}

\answer{城徑二百四十步。}

\working{
立天元一為城徑。以甲東行三百二十步內減二之乙東行三十二步，
餘二百八十八步。以此乘甲東行冪，得二千九百四十九萬一千二百步
為實。一千二百四十八步為廉，四益隅，開三乘方得城徑二百四十步。

\textbf{四乘方程論：}此問所得天元方程為四次方程，中國古代謂之
「三乘方」。西人謂之「四次方程」，以 $x^4$ 為最高次項。
李冶先生之天元術，以籌算排列係數，自下而上為常數項、一次項、
二次項、三次項、四次項，與現代代數學之降冪排列相反，
然其理則一也。開三乘方之法，亦與西洋之數值解法相通，
可見數學為天下之公器，無中外之殊也。
}

\method{
甲東行內減二之乙東行，復以乘甲東行為實，從空，四之甲東行
內減二之乙東行為廉，四益隅，得城徑。
}

%% --- Problem 102 (Vol 6, Problem 17) ---
\problem{
或問：乙從東門直行一十六步而立，甲從乾隅東行三百二十步望乙。
問城徑幾何？又問大差小差各幾何？
}

\answer{城徑二百四十步。大差二百步，小差一百三十六步。}

\working{
立天元一為城徑。甲東行三百二十步減二之乙東行三十二步，
餘二百八十八步。以此乘甲東行三百二十步，得九萬二千一百六十步為實。
一千二百四十八步為廉，四益隅，開三乘方得城徑二百四十步。

求大差：甲東行三百二十步減城徑二百四十步，餘八十步，非大差也。
當以甲東行減半徑：$320 - 120 = 200$ 步為大差之正數。

求小差：乙東行一十六步加半徑一百二十步，得一百三十六步為小差。

\textbf{驗大差小差：}大差二百步，小差一百三十六步。以相乘得
$200 \times 136 = 27200$。大差冪 $= 40000$，小差冪 $= 18496$。
以大小差冪之和減甲東行冪：$102400 - (40000 + 18496) = 43904$，
半之得 $21952$。此數當與大弦之某式相合，以驗天元方程之正確性。
}

\method{
甲東行內減二之乙東行，復以乘甲東行為實，從空，四之甲東行
內減二之乙東行為廉，四益隅，得城徑。
}

%% --- Problem 103 (Vol 6, Problem 18) ---
\problem{
或問：乙從東門直行一十六步，甲從乾隅東行三百二十步望乙，與城參相
直。問城徑幾何？又問通勾股形之弦長幾何？
}

\answer{城徑二百四十步。通弦三百九十二步。}

\working{
立天元一為城徑。甲東行三百二十步減二之乙東行三十二步，
餘二百八十八步。以此乘甲東行冪，得二千九百四十九萬一千二百步為實。
一千二百四十八步為廉，四益隅，開三乘方得城徑二百四十步。

求通弦：通勾股形者，以甲東行為勾、甲南行為股之大勾股形也。
甲東行三百二十步為勾，甲南行六百步為股。以勾股求弦：
$\sqrt{320^2 + 600^2} = \sqrt{102400 + 360000} = \sqrt{462400} = 680$ 步。
然此乃大弦之全數也。其帶分母者，當以大差除之。
大差為甲東行減半徑，即 $320 - 120 = 200$ 步。
以 $680$ 步為分子，$200$ 步為分母，得 $3.4$ 步，非通弦也。
當以天元正術求之：以大差冪加小差冪，減甲東行冪，半之，
加入大小差相乘之積，開方得通弦三百九十二步。

\textbf{按：}通弦之求甚繁，須先求得大小差，然後以大小差之關係
建立方程求之。此問為全卷之終，特設通弦之問，以見天元術之無所不能。
}

\method{
甲東行內減二之乙東行，復以乘甲東行為實，從空，四之甲東行
內減二之乙東行為廉，四益隅，得城徑。
術曰：置甲東行步內減二之乙東行步，餘以乘甲東行冪為實；以四之
甲東行步內減二之乙東行步為廉法；四為益隅法；開三乘方得城徑。
}

\vspace{0.5cm}
\longline
\begin{center}
\textit{測圓海鏡卷六終 --- 凡一十八問}
\end{center}

\newpage

%% ===================================================================
%% APPENDIX: STRUCTURE OF THE WORK
%% ===================================================================
\section*{附錄：全書結構及術語解釋}
\addcontentsline{toc}{section}{附錄：全書結構及術語解釋}

《測圓海鏡》全書十二卷，共一百七十問，為元代李冶先生之傑作。
其書以勾股測圓為主，設圓城一座，甲乙二人從城門出行，以行步之數
推求圓徑。全書結構如下：

\vspace{0.3cm}
\begin{center}
\begin{tabular}{clcc}
\toprule
\textbf{卷} & \textbf{篇名} & \textbf{問數} & \textbf{起止問} \\
\midrule
卷一 & 圓城圖式 & --- & 圖說 \\
卷二 & 正率一十四問 & 14 & 第一問至第十四問 \\
卷三 & 邊股一十七問 & 17 & 第十五問至第三十一問 \\
\rowcolor{gray!15}
卷四 & 底勾一十七問 & 17 & 第三十二問至第四十八問 \\
\rowcolor{gray!15}
卷五 & 大股一十八問 & 18 & 第四十九問至第六十六問 \\
\rowcolor{gray!15}
卷六 & 大勾一十八問 & 18 & 第六十七問至第八十四問 \\
卷七 & 明勾一十六問 & 16 & 第八十五問至第一百問 \\
卷八 & 明股一十六問 & 16 & 第一百一問至第一百二十六問 \\
卷九 & 雜題上 & 18 & 第一百二十七問至第一百四十四問 \\
卷十 & 雜題中 & 18 & 第一百四十五問至第一百六十二問 \\
卷十一 & 雜題下 & 18 & 第一百六十三問至第一百八十問 \\
卷十二 & 雜題終 & 10 & 第一百八十一問至第一百九十問 \\
\midrule
 & \textbf{合計} & \textbf{170} & \\
\bottomrule
\end{tabular}
\end{center}

\vspace{0.5cm}
\noindent\textbf{術語解釋：}

\begin{itemize}
  \item \textbf{天元一}：設未知數為一，即現代代數之 $x$。李冶先生以「立天元一」
  為某數，即設未知數為所求之量。
  
  \item \textbf{寄左}：將某式暫存一邊，相當於方程之一邊。天元術中常以「寄左」
  表示方程之左邊，後以「同數」與之相消。
  
  \item \textbf{相消}：兩式相減，相當於移項合併。天元術之核心步驟，
  以同數與寄左之式相減，得開方式。
  
  \item \textbf{帶分母}：某式含有分母未除。古代算術中常「不受除」，
  即以帶分母之形式保留，以待後續運算。
  
  \item \textbf{冪}：平方，即現代記號 $x^2$。天元術中以某數自乘為其冪。
  
  \item \textbf{開方}：開平方，即求平方根。天元術中最基本之運算。
  
  \item \textbf{開立方}：開立方，即求立方根。用於三次天元方程。
  
  \item \textbf{開三乘方}：開四次方，即求四次根。用於四次天元方程，
  為天元術中最高深之運算。
  
  \item \textbf{實}：常數項，即方程中之已知數部分。
  
  \item \textbf{從}：一次項係數，天元術中「從法」即 $x$ 之係數。
  
  \item \textbf{廉}：二次項係數，即 $x^2$ 之係數。
  
  \item \textbf{隅}：三次項係數（立方方程中），即 $x^3$ 之係數。
  四次方程中另有「三乘隅」或「益隅」之稱。
  
  \item \textbf{勾}：直角三角形之短邊（底邊），相當於現代三角形之底。
  
  \item \textbf{股}：直角三角形之長邊（高邊），相當於現代三角形之高。
  
  \item \textbf{弦}：直角三角形之斜邊，由勾股定理 $\text{勾}^2 + \text{股}^2 = \text{弦}^2$ 求之。
  
  \item \textbf{大勾}：最大勾股形之勾邊，即通勾，為甲從乾隅東行之步數。
  
  \item \textbf{大股}：最大勾股形之股邊，即通股，為甲從乾隅南行之步數。
  
  \item \textbf{大差}：大股減圓徑之餘數，即 $600 - x$ 之類。
  
  \item \textbf{小差}：小勾加半徑之和，或半徑加乙東行之數。
  
  \item \textbf{底勾}：勾股形之最下勾股，與圓徑相切之基本勾股形之勾邊。
\end{itemize}

\vspace{0.5cm}
\noindent\textbf{天元術方程式記法：}

李冶先生以籌算式記多項式，自下而上排列：
\begin{itemize}
  \item 最下為常數項（實）
  \item 其上為一次項（從）
  \item 再其上為二次項（廉）
  \item 再其上為三次項（隅）
\end{itemize}

例如 \tianyuan{480}{1} 表示 $480 - x$（一次式），
\tianyuan{57600}{1} 表示 $57600 - x^2$（二次式），
其中係數為負者以斜線或特殊記號表示之。

\vspace{0.5cm}
\noindent\textbf{各卷方程次數統計：}

\begin{center}
\begin{tabular}{lccc}
\toprule
\textbf{卷} & \textbf{平方（二次）} & \textbf{立方（三次）} & \textbf{三乘方（四次）} \\
\midrule
卷四（底勾） & 12問 & 4問 & 1問 \\
卷五（大股） & 6問 & 10問 & 2問 \\
卷六（大勾） & 3問 & 6問 & 9問 \\
\bottomrule
\end{tabular}
\end{center}

由上表可見，愈至卷六，方程之次數愈高，開方之術愈繁，
乃李冶先生循序漸進、由淺入深之教學法也。

\vspace{1cm}
\begin{center}
\rule{0.3\textwidth}{0.4pt}

\vspace{0.3cm}

{\large\itshape 測圓海鏡卷四至卷六全編終}

\vspace{0.3cm}

{\small 元 \quad 李冶 \quad 撰}

\vspace{0.3cm}

{\small\color{notecolor} 依《欽定四庫全書》子部天文算法類本}
\end{center}

\end{document}
